% ============================================================================
% RELATED: TOPOLOGICAL METHODS AND HIGHER-ORDER STRUCTURE
% ============================================================================
\subsection{Topological Methods and Higher-Order Structure}
\label{sec:related_topological_nn}

\paragraph{Topological Graph Neural Networks.}
Horn et al.\ \citep{horn2022togl} introduced TOGL, a layer that enriches
Graph Neural Networks with global topological information derived from
persistent homology.  Their contribution demonstrates that standard GNN
message-passing, which propagates information along edges, misses
structural invariants visible only at the topological level: connected
components ($H_0$), loops ($H_1$), and voids ($H_2$).  TOGL computes
persistence diagrams from graph filtrations and feeds the resulting
features into the GNN.

The Value architecture addresses structural expressivity through a
different mechanism.  Rather than computing topological summaries from an
existing graph, the system constructs its representations directly in a
high-dimensional binary space where topological properties arise from
the encoding:

\begin{itemize}[nosep]
  \item \textbf{$H_0$ (connected components):} Each Value in a Region
        is a distinct entity.  The degree of overlap between Values---
        measured by their Hamming inner product---defines an implicit
        adjacency structure whose connected components correspond to
        clusters of related inputs.
  \item \textbf{$H_1$ (loops):} If a sequence of boolean operations
        applied to a Value returns it to its original state, this
        constitutes a cycle in the state space.  Detecting such cycles
        is equivalent to identifying $H_1$ features in the
        transformation graph.
\end{itemize}

The trade-off relative to TOGL is computational cost versus
construction cost.  TOGL computes persistence diagrams, which can be
expensive for large graphs, to \emph{detect} topological structure.
The Value architecture does not explicitly detect topology; instead,
any topological structure present in the data is encoded implicitly
through the bit patterns of the resulting Values.

\paragraph{Topological Deep Learning.}
Papillon et al.\ \citep{papillon2024tdl} position topological deep
learning as the next frontier for relational data, arguing that
graph-based representations are limited to pairwise relationships while
real-world data exhibits higher-order interactions best captured by
simplicial complexes, cell complexes, and hypergraphs.

The bitwise OR operation over Values provides a natural encoding of
higher-order relationships: when $n$ Values are combined via disjunction,
the resulting bit pattern captures all $\binom{n}{k}$ subset overlaps
simultaneously through the population-count inner product.  This is not
pairwise attention; it is a set-theoretic encoding where every subset of
the inputs contributes to the composite fingerprint.  The sequential
ordering of operations during Region mixing then imposes a strict
temporal structure on these higher-order relationships.

\paragraph{The Symbolic--Geometric--Symbolic Pipeline.}
A broader trend in recent research involves structured data being mapped
to a vector space, processed via geometric operations, and reconstructed
as interpretable symbols.  The Value architecture implements a
compressed version of this pipeline: raw bytes are deterministically
mapped to binary frames, processed through boolean operations during
Region mixing, and the resulting frames are decoded back to byte
sequences.  The entire pipeline is deterministic and executes without
gradient descent.
